% chapters/03/section-04.tex
\section{Struktur Folder Proyek dan Pengelolaan Versi}
\label{sec:struktur-git}

Sebelum menulis program, tempat bekerja perlu disiapkan. Bagian ini menjelaskan
struktur folder yang dipakai sepanjang semester dan cara menyimpan pekerjaan
dengan Git.

\subsection{Struktur Folder Proyek}

Struktur folder yang dipakai diperlihatkan pada Gambar~\ref{fig:struktur-proyek}.
Susunannya sengaja dibuat sederhana, karena yang penting bukan kerumitannya,
melainkan pemisahan berkas menurut fungsinya.

\begin{figure}[htbp]
\centering
% chapters/03/figures/struktur-proyek.tikz
% Struktur folder proyek yang dipakai sepanjang semester.
% Pohon digambar mendatar agar nama berkas yang panjang tetap muat.
\begin{tikzpicture}[
  dir/.style={
    rectangle, rounded corners=1pt, draw=ubPrimary, thick, fill=ubAlt,
    inner sep=3.5pt, font=\ttfamily\footnotesize, text=ubPrimary
  },
  berkas/.style={
    rectangle, rounded corners=1pt, draw=ubLine, fill=white,
    inner sep=3.5pt, font=\ttfamily\footnotesize
  },
  grow'=right,
  level 1/.style={sibling distance=9mm, level distance=40mm},
  level 2/.style={sibling distance=7mm, level distance=40mm},
  edge from parent/.style={draw=ubAccent, thick}
]

\node[dir, anchor=west] (akar) at (-6.2,0) {proyek-pygame}
  child { node[dir] {code}
    child { node[berkas] {ch03\_jendela.py} }
    child { node[berkas] {ch05\_gerak.py} } }
  child { node[dir] {catatan}
    child { node[berkas] {praktikum-01.md} } }
  child { node[dir] {aset} }
  child { node[berkas] {.gitignore} }
  child { node[berkas] {README.md} };

\node[font=\scriptsize\itshape, ubPrimary, align=center, text width=112mm]
  at (0,-3.3)
  {Nama folder dan berkas ditulis dengan huruf kecil dan garis bawah, tanpa
   spasi, agar aman dipakai pada semua sistem operasi.};
\end{tikzpicture}

\caption{Struktur folder proyek yang dipakai sepanjang semester}
\label{fig:struktur-proyek}
\end{figure}

Tiga folder utama memiliki tugas masing-masing.

\begin{description}
  \item[\code{code}] Berisi seluruh berkas program, satu berkas untuk setiap
        latihan atau praktikum. Pemisahan ini memudahkan Anda kembali ke
        program tertentu ketika diperlukan.
  \item[\code{catatan}] Berisi catatan praktikum: apa yang dikerjakan, kendala
        yang muncul, dan cara penyelesaiannya. Catatan inilah yang menjadi
        bahan laporan pada proyek akhir nanti.
  \item[\code{aset}] Berisi berkas gambar dan bunyi yang dipakai program.
        Folder ini baru terisi mulai Bab~VII, ketika program mulai memuat
        gambar.
\end{description}

Dua berkas tambahan melengkapinya. Berkas \code{README.md} menerangkan isi
repositori, dan berkas \code{.gitignore} menentukan berkas mana yang tidak
perlu ikut disimpan.

Perhatikan penamaan folder dan berkas: seluruhnya huruf kecil, memakai garis
bawah sebagai pengganti spasi, dan tanpa karakter khusus. Kebiasaan itu
menghindari masalah ketika proyek dibuka pada sistem operasi lain atau
dijalankan melalui terminal, tempat spasi pada nama berkas harus ditulis
dengan cara khusus.

\subsection{Berkas \code{.gitignore}}

Tidak semua berkas perlu disimpan. Ketika program Python dijalankan, Python
membuat folder \code{__pycache__} berisi salinan sementara yang mempercepat
pemuatan berikutnya. Berkas itu dapat dibuat ulang kapan saja dan berbeda pada
setiap komputer, sehingga menyimpannya hanya membuat repositori membengkak.

Buatlah berkas bernama \code{.gitignore} pada folder teratas proyek dan isi
dengan baris berikut.

\begin{kodekeluaran}
__pycache__/
*.pyc
.venv/
.vscode/
\end{kodekeluaran}

Baris pertama mengabaikan folder salinan sementara Python. Baris kedua
mengabaikan berkas terjemahan kode. Baris ketiga mengabaikan folder lingkungan
virtual, dan baris keempat mengabaikan pengaturan editor yang bersifat
pribadi.

\subsection{Memulai Repositori Git}

Git bekerja dengan tiga tempat penyimpanan yang perlu Anda bedakan sejak awal.

\begin{enumerate}
  \item \textbf{Folder kerja} (\istilah{working directory}), yaitu berkas yang
        sedang Anda sunting.
  \item \textbf{Daftar siap simpan} (\istilah{staging area}), yaitu berkas
        yang sudah ditandai untuk ikut disimpan pada penyimpanan berikutnya.
  \item \textbf{Riwayat} (\istilah{repository}), yaitu kumpulan penyimpanan
        yang telah tercatat dan tidak berubah lagi.
\end{enumerate}

Alur kerjanya berpindah dari folder kerja ke daftar siap simpan, lalu ke
riwayat. Perintah untuk itu dijalankan berurutan sebagai berikut.

\begin{kodekeluaran}
cd proyek-pygame          (masuk ke folder proyek)
git init                  (sekali saja: memulai repositori)
git status                (memeriksa keadaan berkas)

git add .                 (menandai seluruh perubahan)
git commit -m "praktikum 01: jendela pertama dan struktur proyek"
git log --oneline         (melihat riwayat penyimpanan)
\end{kodekeluaran}

Perintah \code{git add} tidak menyimpan apa pun ke riwayat; ia hanya
memindahkan perubahan ke daftar siap simpan. Penyimpanan sebenarnya terjadi
pada perintah \code{git commit}. Pemisahan dua langkah ini bermanfaat ketika
Anda mengubah banyak berkas sekaligus: Anda dapat memilih hanya berkas yang
berkaitan dengan satu pekerjaan, sehingga riwayatnya tetap mudah dibaca.

\begin{tipsbox}
Beberapa anjuran yang terbukti menghemat waktu:
\begin{itemize}
  \item Simpanlah sesering mungkin, misalnya setiap kali satu langkah
        praktikum selesai. Riwayat yang rinci memudahkan Anda kembali ke
        keadaan terakhir yang berjalan baik.
  \item Tulis pesan penyimpanan yang menerangkan isinya, bukan sekadar
        ``update'' atau ``perbaikan''. Sebutkan apa yang dikerjakan.
  \item Periksa hasil dengan \code{git status} sebelum dan sesudah menyimpan,
        terutama pada minggu-minggu pertama, agar Anda mengetahui berkas mana
        yang sudah tercatat.
  \item Apabila sebuah percobaan gagal dan Anda ingin kembali, jangan menghapus
        berkas secara manual. Simpan keadaan gagal itu, lalu perbaiki; riwayat
        yang jujur justru berguna sebagai catatan belajar.
\end{itemize}
\end{tipsbox}

\subsection{Repositori sebagai Bukti Portofolio}

Sepuluh persen dari nilai mata kuliah ini berasal dari penilaian artefak kode
mingguan, dan dosen menilai enam sampel secara acak. Konsekuensinya, yang
dinilai bukan hanya program terakhir pada minggu kelima belas, melainkan
kesinambungan pekerjaan Anda.

Repositori yang tertata rapi juga membantu Anda sendiri. Ketika pada Bab~X
Anda perlu mengingat kembali bagaimana gerak objek dihitung pada Bab~V,
riwayat pekerjaan Anda adalah tempat pertama yang sebaiknya dibuka.
